\chapter{信号処理の基礎}


\section{この章の目的}

本章では，信号処理を学び始めるときに必要になる最小限の概念を整理します．
ここでいう\textbf{信号}とは，時間，空間，周波数などの軸に沿って並んだ数値列です．
音声波形，画像の画素値列，センサの時系列，ニューラルネットワークに入力する特徴量列は，いずれも信号として扱えます．

機械学習の前処理では，生の観測値をそのままモデルに入れるのではなく，
\begin{itemize}
\item どの軸に沿って値が並んでいるか
\item どのくらい細かく観測されているか
\item どのような操作で情報が強調または失われるか
\end{itemize}
を理解しておく必要があります．

本章の到達目標は次の通りです．
\begin{itemize}
\item 離散時間信号を数列として表現できる
\item サンプリング周期とサンプリング周波数の意味を説明できる
\item インパルス，ステップ，正弦波といった基本信号を説明できる
\item シフト，反転，スケーリングなどの基本操作が何をするか理解できる
\item エネルギーと平均電力の違いを説明できる
\end{itemize}


\section{信号とは何か}

連続時間で変化する物理量を $t$ の関数 $x(t)$ と表すことがあります．
例えばマイクで観測される空気圧の変化は連続時間信号として考えられます．
一方，コンピュータ上では通常，信号は一定間隔で標本化された数列として扱います．
このとき離散時間信号を
\[
x[n], \quad n = 0, 1, 2, \ldots
\]
と書きます．

$n$ は\textbf{サンプル番号}であり，実際の時間そのものではありません．
サンプリング周期を $T_s$ 秒とすると，$n$ 番目のサンプルは実時間では
\[
t = nT_s
\]
に対応します．
したがって，離散時間信号は「時間の情報を失った数列」ではなく，
\textbf{一定間隔で観測した連続時間信号の近似表現}です．


\section{サンプリング周期とサンプリング周波数}

サンプリング周期 $T_s$ は，隣り合う 2 つのサンプルの時間間隔です．
その逆数
\[
f_s = \frac{1}{T_s}
\]
を\textbf{サンプリング周波数}と呼びます．単位は Hz です．

例えば $f_s = 16000$ Hz なら，1 秒間に 16000 点観測していることになります．
音声処理では 16 kHz や 44.1 kHz，画像処理では画素間隔やフレームレートが同様の役割を持ちます．

ここで重要なのは，離散信号だけを見ても「隣り合うサンプルが何秒離れているか」は自動では分からないことです．
同じ数列でも，$f_s$ が変われば意味する物理現象は変わります．
したがって，信号データを扱うときは，値そのものだけでなく\textbf{サンプリング条件も一緒に管理する}必要があります．


\section{基本信号}

\subsection{単位インパルス}

\textbf{単位インパルス}は
\[
\delta[n] =
\begin{cases}
1 & (n = 0) \\
0 & (n \neq 0)
\end{cases}
\]
で定義されます．

一見すると単純な信号ですが，インパルスはシステムの性質を調べるための基準信号として非常に重要です．
後の LTI 系の章では，インパルス応答から系全体の振る舞いを記述できることを学びます．


\subsection{単位ステップ}

\textbf{単位ステップ}は
\[
u[n] =
\begin{cases}
1 & (n \geq 0) \\
0 & (n < 0)
\end{cases}
\]
で定義されます．

信号がある時刻から立ち上がる状況や，累積和のような振る舞いを考えるときに役立ちます．


\subsection{正弦波}

離散時間の正弦波は
\[
x[n] = A \sin(\omega_0 n + \phi)
\]
と書けます．

\begin{itemize}
\item $A$ は振幅
\item $\omega_0$ は離散時間角周波数
\item $\phi$ は初期位相
\end{itemize}

です．

連続時間での周波数 $f$ Hz とサンプリング周波数 $f_s$ Hz の関係は
\[
\omega_0 = 2\pi \frac{f}{f_s}
\]
です．
この式から，離散時間での正弦波の見え方は，元の周波数 $f$ だけでなくサンプリング周波数 $f_s$ にも依存することが分かります．


\section{基本操作}

\subsection{時間シフト}

信号を $n_0$ サンプルだけ遅らせる操作は
\[
y[n] = x[n - n_0]
\]
です．
逆に $x[n+n_0]$ は $n_0$ サンプルだけ進めた信号です．

離散時間では，シフト量は通常整数サンプルで扱います．
音響では遅延，画像では平行移動に対応する基本操作です．


\subsection{反転}

\[
y[n] = x[-n]
\]
は時間反転です．
畳み込みを理解する際に必ず現れる操作なので，単なる見た目の変換としてではなく，後で出てくる演算の一部として覚える必要があります．


\subsection{振幅スケーリング}

\[
y[n] = a x[n]
\]
は振幅を $a$ 倍する操作です．
$|a| > 1$ なら増幅，$|a| < 1$ なら減衰です．
$a < 0$ なら上下反転も同時に起こります．


\section{エネルギーと平均電力}

信号の大きさを要約する尺度として，\textbf{エネルギー}と\textbf{平均電力}があります．

離散時間信号 $x[n]$ のエネルギーを
\[
E = \sum_{n=-\infty}^{\infty} |x[n]|^2
\]
と定義します．

有限長の信号や，時間とともに十分小さくなる信号ではエネルギーが有限になります．
一方で，ずっと振動し続ける正弦波のような信号は，無限に長く観測するとエネルギーが発散します．

そのような信号には平均電力
\[
P = \lim_{N \to \infty} \frac{1}{2N+1}\sum_{n=-N}^{N} |x[n]|^2
\]
を使います．

有限長パルスはエネルギー信号，定常的に続く正弦波は電力信号として扱うのが典型です．
どちらを使うべきかは，対象が「有限のイベント」なのか「定常的に続く現象」なのかで決まります．


\section{Python で基本信号を確かめる}

以下のコードは，正弦波とノイズを生成し，サンプル番号と実時間の両方を意識して可視化する最小例です．

\begin{lstlisting}[style=codecell,language=Python]
from pathlib import Path

import matplotlib.pyplot as plt
import numpy as np

output_dir = Path("outputs/ch13")
output_dir.mkdir(parents=True, exist_ok=True)

fs = 16000
duration = 0.01
n = np.arange(int(fs * duration))
t = n / fs

clean = np.sin(2 * np.pi * 440 * t)
rng = np.random.default_rng(0)
noise = 0.1 * rng.standard_normal(size=n.shape)
observed = clean + noise

plt.figure(figsize=(8, 3))
plt.plot(n, observed, label="observed", linewidth=1)
plt.plot(n, clean, label="clean", linewidth=2)
plt.xlabel("sample index n")
plt.ylabel("amplitude")
plt.legend()
plt.tight_layout()
plt.savefig(output_dir / "signal_and_noise.png", dpi=150)
plt.close()
\end{lstlisting}

このコードで確認したい点は 2 つあります．
\begin{itemize}
\item $n$ を横軸にすると離散信号は数列として見える
\item $t=n/f_s$ を使えば，同じ信号を実時間軸でも解釈できる
\end{itemize}

モデルに入力する直前のデータが「何番目のサンプルか」だけでなく「何秒に対応するか」まで追えるようになると，
フィルタ設計，窓長の設定，特徴量抽出の妥当性を判断しやすくなります．


\section{この章で押さえるべき点}

\begin{itemize}
\item 離散時間信号 $x[n]$ は，サンプル番号に対応づけられた数列である
\item サンプリング周波数 $f_s$ がなければ，信号の時間的な意味は確定しない
\item インパルス，ステップ，正弦波は後の解析の基準になる基本信号である
\item シフト，反転，スケーリングは畳み込みやフィルタ理解の前提になる
\item 有限長信号と定常信号では，エネルギーと平均電力のどちらを見るべきかが異なる
\end{itemize}

次章では，正弦波を基底として信号を周波数ごとに分解する考え方，すなわちスペクトル解析を導入します．
